\chapter{Platelet Count}

\section*{What Is This Test?}
The platelet count, also called the thrombocyte count, measures the number of platelets (thrombocytes) per unit volume of whole blood. Platelets are small (2--4 $\mu$m diameter), anucleate cytoplasmic fragments derived from megakaryocytes in the bone marrow under the regulation of thrombopoietin. They circulate in the blood for approximately 7--10 days before being removed by the reticuloendothelial system, primarily in the spleen. Platelets play a critical role in primary hemostasis: when a blood vessel is injured, platelets adhere to exposed subendothelial collagen (via von Willebrand factor and glycoprotein Ib), activate, change shape, release granular contents (ADP, thromboxane A2, serotonin), and aggregate to form a platelet plug that seals the site of vascular injury. The platelet count is a core component of the complete blood count and is essential for evaluating bleeding disorders, thrombotic conditions, and bone marrow function. Thrombocytopenia (decreased platelet count) is the most common cause of clinically significant bleeding, while thrombocytosis (elevated platelet count) can be reactive (secondary to inflammation, infection, or iron deficiency) or primary (resulting from myeloproliferative neoplasms such as essential thrombocythemia) \cite{tierney2023current}.

\section*{Alternative Names}
Plt; Thrombocyte Count; Platelet Number; PLT Count; Thrombocyte Concentration

\section*{Principle}
Platelets are counted on automated hematology analyzers using impedance (Coulter) technology, optical light scatter, or fluorescent flow cytometry. In impedance counting, whole blood is diluted and red blood cells are lysed. The remaining platelets are passed through a small aperture; each platelet generates an electrical resistance pulse proportional to its size. Platelets are discriminated from other particles by their characteristic size range (2--20 fL). The analyzer counts 5,000--10,000 platelets per sample, giving a coefficient of variation of approximately 3--5\% \cite{tietz2012textbook}. Optical platelet counting (Sysmex XN, Siemens ADVIA) uses laser light scatter and/or fluorescent nucleic acid dyes (such as polymethine or oxazine) that specifically label platelets, providing more accurate counts in severely thrombocytopenic samples. The immature platelet fraction (IPF), analogous to the reticulocyte count for red cells, measures young, RNA-rich platelets newly released from megakaryocytes and is an indicator of bone marrow thrombopoietic activity \cite{kaushansky2021williams}. Manual platelet counting using phase-contrast microscopy and a Neubauer hemocytometer chamber (diluted with 1\% ammonium oxalate to lyse red blood cells) is a reference method, rarely used clinically, but useful when automated counts are flagged as unreliable (e.g., severe thrombocytopenia, giant platelets, platelet clumping) \cite{bain2017dacie}. Platelet counts from EDTA-anticoagulated blood must be verified when pseudothrombocytopenia is suspected (platelet clumping in EDTA due to antiplatelet antibodies); recollecting in sodium citrate tube and correcting for dilution is standard practice.

\section*{History}
Platelets were first described in the 1840s-1880s by several observers who noted small ``corpuscles'' in blood, but their role in hemostasis was not recognized until Giulio Bizzozero in 1882 demonstrated that they were the first cellular element to adhere to sites of vascular injury and named them ``piastrine'' (little plates). James Homer Wright, in 1906, established that platelets originate from bone marrow megakaryocytes. The first practical method for counting platelets manually using a hemocytometer and phase-contrast microscopy was developed by Brecher and Cronkite in 1950. The Coulter principle, patented in 1953 and commercialized in 1956, automated platelet counting with far greater precision and speed. The introduction of flow cytometric methods and the immature platelet fraction in the early 2000s further enhanced platelet assessment. Today, automated platelet counts are among the most commonly ordered laboratory tests worldwide.

\section*{Why This Test Is Ordered}
\begin{itemize}
  \item Evaluation of bleeding disorders --- thrombocytopenia is the most common cause of mucocutaneous bleeding (petechiae, purpura, epistaxis, gingival bleeding, menorrhagia); spontaneous bleeding risk increases when platelets fall below 50,000/$\mu$L and becomes severe below 10,000/$\mu$L
  \item Screening before surgery or invasive procedures --- a platelet count >50,000/$\mu$L is generally adequate for most procedures; >80,000--100,000/$\mu$L for neurosurgery or ophthalmic surgery
  \item Diagnosis and monitoring of immune thrombocytopenic purpura (ITP) --- isolated thrombocytopenia with normal WBC and hemoglobin, large/giant platelets on smear, and exclusion of other causes
  \item Diagnosis of thrombotic thrombocytopenic purpura (TTP) --- severe thrombocytopenia with microangiopathic hemolytic anemia (schistocytes), fever, neurologic symptoms, and renal dysfunction; requires emergent plasma exchange
  \item Evaluation of disseminated intravascular coagulation (DIC) --- consumptive thrombocytopenia with prolonged PT/PTT, elevated D-dimer, low fibrinogen, and schistocytes
  \item Diagnosis of heparin-induced thrombocytopenia (HIT) --- thrombocytopenia occurring 5--14 days after heparin exposure, with thrombotic complications; confirmed by anti-PF4/heparin antibody ELISA and serotonin release assay
  \item Monitoring of chemotherapy and radiation therapy --- expected thrombocytopenia with nadir 10--14 days after chemotherapy
  \item Diagnosis of essential thrombocythemia and other myeloproliferative neoplasms --- persistent, unexplained thrombocytosis >450,000/$\mu$L with JAK2 V617F, CALR, or MPL mutation
  \item Preoperative screening and evaluation of liver disease (thrombocytopenia from hypersplenism in cirrhosis with portal hypertension)
  \item Monitoring of heparin therapy and other anticoagulants
  \item Evaluation of splenomegaly and hypersplenism
\end{itemize}

\section*{Primary vs Secondary Test}
The platelet count is a primary component of the complete blood count and is essential for the evaluation of bleeding disorders, thrombotic conditions, and bone marrow function.

\section*{How This Test Is Performed}
A venous blood sample is collected in a lavender-top (EDTA) tube from a peripheral vein, typically the antecubital fossa. Approximately 2--5 mL of blood is drawn, and the tube is gently inverted 8--10 times to ensure adequate mixing with the EDTA anticoagulant. The sample is processed on an automated hematology analyzer, which counts platelets by impedance, optical, or fluorescence methods \cite{fischbach2023manual}. The analysis is completed in 1--2 minutes. If the automated platelet count is flagged for platelet clumps, giant platelets, or very low platelet count (<30,000/$\mu$L), the peripheral blood smear is examined to estimate platelet count manually and confirm or exclude pseudothrombocytopenia. If pseudothrombocytopenia (EDTA-induced platelet clumping) is suspected, a fresh sample is collected in a light blue-top (sodium citrate) tube, and the obtained platelet count is multiplied by 1.1 to correct for the dilution effect of liquid citrate anticoagulant. The venipuncture procedure itself takes less than 5 minutes.

\section*{What Preparation Is Needed}
No special preparation is required. Fasting is not necessary. The patient should inform the healthcare provider of all medications, with particular attention to antiplatelet agents (aspirin, clopidogrel, ticagrelor), anticoagulants (heparin, warfarin, direct oral anticoagulants), chemotherapy, immunosuppressants, antibiotics (vancomycin, linezolid), anticonvulsants (valproic acid, phenytoin), and over-the-counter supplements that can affect platelet function (fish oil, vitamin E, ginkgo biloba, garlic). Recent blood transfusion, surgery, or trauma should be disclosed. Heavy alcohol consumption can cause thrombocytopenia through direct marrow toxicity.

\section*{Contraindications and Precautions}
No absolute contraindications to blood collection. Critical pre-analytical and analytical considerations include: EDTA is the standard anticoagulant, but EDTA-induced pseudothrombocytopenia occurs in approximately 0.1--0.2\% of individuals due to EDTA-dependent antiplatelet antibodies causing platelet clumping at room temperature --- this is a laboratory artifact (not a clinical disorder) and must be recognized to avoid unnecessary investigation and intervention \cite{wallach2022interpretation}; the sample must be thoroughly mixed immediately after collection to prevent microclots that consume platelets and falsely lower the count; hemolyzed or clotted specimens are unacceptable; prolonged EDTA storage (>4 hours at room temperature) causes platelet swelling and disintegration, decreasing the count; very high platelet counts (>1,000,000/$\mu$L) may exceed the analyzer's linear range and require dilution; giant platelets (as in ITP, Bernard-Soulier syndrome, May-Hegglin anomaly) may be excluded from the platelet gate by impedance analyzers, yielding falsely low counts --- optical or fluorescence methods, or manual smear estimation, are advised; platelet satellitism (platelets adhering to neutrophils in EDTA) causes a falsely low count; lipemia, cryoglobulins, and red cell fragments (schistocytes in TTP/HUS) may interfere with automated counting. Samples should be analyzed within 4 hours at room temperature or 24 hours if refrigerated.

\section*{Risks}
Minimal risks of venipuncture: mild pain or stinging, bruising (ecchymosis), hematoma, lightheadedness or vasovagal syncope, and extremely rarely, infection or nerve injury. In patients with severe thrombocytopenia, prolonged pressure should be applied to the venipuncture site to prevent hematoma. The use of a small-gauge needle (21--23 gauge) is recommended.

\section*{What Happens After the Test}
Platelet count results are available within 30--60 minutes as part of the CBC. Critical values requiring immediate notification typically include platelet count <20,000/$\mu$L (risk of spontaneous intracranial or gastrointestinal bleeding) or >1,000,000/$\mu$L (risk of thrombosis or bleeding due to acquired von Willebrand disease in extreme thrombocytosis). The platelet count is interpreted in the context of the full CBC, peripheral blood smear morphology, and clinical presentation. For thrombocytopenia, the smear is reviewed for platelet clumps (pseudothrombocytopenia), schistocytes (TTP/HUS/DIC), and blast cells (leukemia). For thrombocytosis, iron studies are checked for iron deficiency anemia, and JAK2/CALR/MPL mutation testing is considered for suspected essential thrombocythemia. Serial platelet counts are monitored during chemotherapy, heparin therapy, and recovery from bone marrow transplantation.

\section*{Understanding Results}

\subsection*{Below Normal (Thrombocytopenia: <150,000/$\mu$L)}
\begin{itemize}
  \item \textbf{Immune thrombocytopenic purpura (ITP):} Isolated thrombocytopenia (platelets often 10,000--50,000/$\mu$L) with otherwise normal CBC and peripheral smear. Large/giant platelets present due to accelerated thrombopoiesis. Caused by autoantibodies (typically IgG) against platelet glycoproteins (GP IIb/IIIa, GP Ib/IX). Diagnosed by exclusion after ruling out drug-induced, infection-related, and other secondary causes. Bone marrow examination shows normal or increased megakaryocytes \cite{neunert2019american}.
  \item \textbf{Drug-induced thrombocytopenia:} Immune-mediated (drug-dependent antibodies) from heparin, quinine, quinidine, sulfonamides, vancomycin, rifampin, carbamazepine, valproic acid, and many others. Non-immune (dose-dependent marrow suppression) from chemotherapy, alcohol (direct megakaryocyte toxicity), thiazide diuretics, and valproic acid.
  \item \textbf{Heparin-induced thrombocytopenia (HIT):} Immune-mediated type II HIT occurs 5--14 days after heparin exposure (earlier if recent prior heparin exposure). Platelets typically fall 30--50\% from baseline to 50,000--80,000/$\mu$L. Paradoxically associated with thrombosis (venous and arterial). Caused by IgG antibodies against heparin-PF4 complexes that activate platelets. Diagnosed by clinical scoring (4Ts score), anti-PF4/heparin ELISA, and serotonin release assay. Requires immediate discontinuation of all heparin and initiation of non-heparin anticoagulant (argatroban, bivalirudin, fondaparinux) \cite{greinacher2015heparin}.
  \item \textbf{Thrombotic thrombocytopenic purpura (TTP):} Severe thrombocytopenia (often <20,000/$\mu$L) with microangiopathic hemolytic anemia (schistocytes on smear, elevated LDH, low haptoglobin), fever, fluctuating neurologic symptoms, and renal dysfunction. Caused by deficiency of ADAMTS13 (acquired autoantibody or congenital), leading to uncleaved ultra-large von Willebrand factor multimers causing platelet microthrombi. Mortality approaches 90\% without treatment; requires emergent plasma exchange and corticosteroids.
  \item \textbf{Disseminated intravascular coagulation (DIC):} Consumptive thrombocytopenia with prolonged PT and PTT, elevated D-dimer, low fibrinogen, and schistocytes. Caused by sepsis, trauma, obstetric emergencies, malignancy, and severe pancreatitis.
  \item \textbf{Sepsis:} One of the most common causes of thrombocytopenia in hospitalized patients, through multifactorial mechanisms including DIC, immune destruction, and marrow suppression.
  \item \textbf{Liver disease and hypersplenism:} Splenic sequestration of platelets in cirrhosis with portal hypertension and splenomegaly; platelet count is typically mild to moderately reduced (50,000--100,000/$\mu$L).
  \item \textbf{Bone marrow failure:} Aplastic anemia, myelodysplastic syndrome, leukemia, lymphoma, multiple myeloma, metastatic carcinoma, myelofibrosis, and granulomatous infiltration.
  \item \textbf{Chemotherapy and radiation therapy:} Dose-dependent myelosuppression with nadir at 10--14 days.
  \item \textbf{Viral infections:} HIV, Epstein-Barr virus, cytomegalovirus, hepatitis C, dengue fever, parvovirus B19.
  \item \textbf{Pregnancy:} Gestational thrombocytopenia (benign, mild, 70,000--150,000/$\mu$L, occurs in 5--8\% of pregnancies, third trimester); preeclampsia/HELLP syndrome (hemolysis, elevated liver enzymes, low platelets); acute fatty liver of pregnancy.
  \item \textbf{Post-transfusion purpura:} Severe, sudden thrombocytopenia 5--10 days after blood transfusion, caused by alloantibodies against platelet antigens (especially HPA-1a).
  \item \textbf{Thrombocytopenia in massive transfusion:} Dilutional thrombocytopenia when >1.5 blood volumes are replaced with packed red blood cells without platelet transfusion.
  \item \textbf{Congenital thrombocytopenia:} Bernard-Soulier syndrome (giant platelets, GP Ib deficiency), Wiskott-Aldrich syndrome (microthrombocytopenia, immunodeficiency, eczema), May-Hegglin anomaly (giant platelets with D\"{o}hle bodies in neutrophils), thrombocytopenia-absent radius (TAR) syndrome.
\end{itemize}

\subsection*{Above Normal (Thrombocytosis: >450,000/$\mu$L)}
\begin{itemize}
  \item \textbf{Reactive (secondary) thrombocytosis:} The most common cause of elevated platelet count. Caused by acute infection, chronic inflammation (rheumatoid arthritis, inflammatory bowel disease), iron deficiency anemia (platelets often 500,000--800,000/$\mu$L that normalize with iron replacement), acute blood loss, hemolytic anemia, post-splenectomy (loss of splenic platelet sequestration causes mild persistent thrombocytosis 500,000--1,000,000/$\mu$L), malignancy (paraneoplastic), recovery from thrombocytopenia (rebound), and vigorous exercise. Typically platelets <1,000,000/$\mu$L, normal platelet morphology and function, no increased thrombotic risk.
  \item \textbf{Essential thrombocythemia (ET):} Primary myeloproliferative neoplasm characterized by persistent, unexplained thrombocytosis >450,000/$\mu$L (often >1,000,000/$\mu$L). Caused by clonal proliferation of megakaryocytes. Mutations: JAK2 V617F (50--60\%), CALR (15--25\%), MPL (5\%), and triple-negative (10--15\%) \cite{barbui2018who}. Giant platelets and megakaryocytic fragments may be seen on smear. Associated with thrombotic complications (cerebral, coronary, peripheral arterial thrombosis; deep vein thrombosis; portal/hepatic vein thrombosis in younger patients), erythromelalgia, and paradoxically, bleeding (especially with extreme thrombocytosis >1,500,000/$\mu$L due to acquired von Willebrand disease from adsorption of large vWF multimers onto platelets).
  \item \textbf{Polycythemia vera:} Elevated platelet count in addition to elevated hemoglobin and hematocrit; JAK2 V617F mutation.
  \item \textbf{Primary myelofibrosis:} Thrombocytosis often accompanies leukoerythroblastic blood picture with teardrop cells.
  \item \textbf{Chronic myeloid leukemia:} Thrombocytosis may be present along with marked leukocytosis and basophilia; Philadelphia chromosome and BCR-ABL1 confirm CML.
  \item \textbf{Iron deficiency anemia:} Mild to moderate thrombocytosis (500,000--800,000/$\mu$L) that resolves with iron therapy.
\end{itemize}

\section*{Normal Results}
\begin{itemize}
  \item \textbf{Adults:} 150,000--400,000/$\mu$L (150--400 $\times$ 10\textsuperscript{9}/L)
  \item \textbf{Children (1--12 years):} 150,000--450,000/$\mu$L
  \item \textbf{Infants (1--12 months):} 200,000--450,000/$\mu$L
  \item \textbf{Newborns (first week):} 150,000--400,000/$\mu$L (may be mildly lower in preterm infants)
  \item \textbf{Pregnancy (gestational thrombocytopenia):} Mild decrease, with normal range 110,000--350,000/$\mu$L; counts <70,000/$\mu$L require investigation for ITP, preeclampsia, or other causes
  \item \textbf{Post-splenectomy:} 500,000--1,000,000/$\mu$L (expected postoperative thrombocytosis)
  \item \textbf{High altitude:} Mild increase due to compensatory erythrocytosis and increased thrombopoietin
\end{itemize}

\section*{Follow-up Tests}
\begin{itemize}
  \item \textbf{Peripheral blood smear:} Essential for assessing platelet size (giant, small), platelet clumping (pseudothrombocytopenia), schistocytes (TTP/HUS/DIC), blast cells (leukemia), dysplastic features, and white/red cell morphology
  \item \textbf{Repeat platelet count in sodium citrate tube:} When EDTA-induced pseudothrombocytopenia is suspected (platelet clumps on smear)
  \item \textbf{Complete blood count with WBC differential:} Evaluate for concurrent cytopenias or abnormalities
  \item \textbf{Bone marrow aspiration and biopsy:} When pancytopenia is present (aplastic anemia), suspected leukemia, MDS, or myelofibrosis
  \item \textbf{ADAMTS13 activity and inhibitor:} When TTP is suspected (severe thrombocytopenia with microangiopathic hemolytic anemia and neurologic/renal involvement)
  \item \textbf{Coagulation panel (PT, PTT, fibrinogen, D-dimer):} When DIC is suspected
  \item \textbf{Anti-PF4/heparin antibody ELISA and serotonin release assay:} When HIT is suspected (thrombocytopenia 5--14 days after heparin exposure)
  \item \textbf{HIV, hepatitis C, EBV, CMV serologies:} For infection-associated thrombocytopenia
  \item \textbf{JAK2 V617F, CALR, and MPL mutation testing:} When essential thrombocythemia or other myeloproliferative neoplasm is suspected with persistent thrombocytosis
  \item \textbf{Antiplatelet antibody testing:} Limited utility for ITP diagnosis (low sensitivity/specificity); diagnosis is clinical by exclusion
  \item \textbf{Iron studies (ferritin, serum iron, TIBC):} When iron deficiency anemia is suspected with thrombocytosis
  \item \textbf{Liver function tests and abdominal ultrasound:} When liver disease and hypersplenism are suspected
  \item \textbf{Flow cytometry:} When leukemia or lymphoma is suspected
  \item \textbf{von Willebrand factor panel (vWF antigen, ristocetin cofactor activity, factor VIII):} When extreme thrombocytosis (>1,500,000/$\mu$L) or bleeding is present, to assess for acquired von Willebrand disease
\end{itemize}

\section*{References}
\bibliographystyle{unsrtnat}
\bibliography{references}

\clearpage
\section*{Regulatory Status and Authorization Date}
Regulatory status applies to the specific assay, instrument, manufacturer, and intended use rather than to the test name alone. No single approval date can be assigned to this test category. For a commercial assay, verify the current product in the relevant regulator's database: the U.S. Food and Drug Administration (FDA) clearance, approval, or authorization; India's Central Drugs Standard Control Organisation (CDSCO) licence or permission; the European Union In Vitro Diagnostic Regulation (EU IVDR) CE certificate issued through the applicable conformity-assessment route; or the United Kingdom Medicines and Healthcare products Regulatory Agency (MHRA) registration and UKCA/CE status. Laboratory-developed tests may not have an individual FDA, CDSCO, EU IVDR, or MHRA approval date.
\clearpage
